\documentclass[onecolumn,12pt]{IEEEtran}
\IEEEoverridecommandlockouts

% Essential packages
\usepackage{cite}
\usepackage{amsmath,amssymb,amsfonts}
\usepackage{graphicx}
\usepackage{xcolor}
\usepackage{caption}
\usepackage{fancyhdr}
\usepackage{url}
\usepackage{enumitem}
\usepackage{subcaption}

\usepackage{microtype}

% Caption formatting
\captionsetup{font=footnotesize}

% Header setup
\pagestyle{fancy}
\fancyhf{}
\fancyhead[C]{\textbf{Research Proposal}\\\rule{\textwidth}{0.4pt}}
\renewcommand{\headrulewidth}{0pt}

\begin{document}

\title{Efficient 9-Qubit NEQR RGB Patch Encoding and Trainable Quantum Filters for Hybrid Quantum-Classical ResNet-18}
\author{
University of Nottingham\\
School of Computer Science\\
\small{Professor Dr. Tan Chye Cheah}\\
\small{Joseph Emmanuel Chay Huan Qin}
}
\maketitle

\begin{abstract}
This work overcomes the lack of full‐color, multi‐qubit patch encoders in quantum convolutional neural networks by introducing a 9‐qubit NEQR amplitude encoder for $3\times3\times3$ RGB patches, combined with Gray‐code amplitude loading and quantum‐PCA compression to reduce state‐preparation depth by up to 40\%. We design trainable variational filters using both hardware‐efficient and tree‐tensor‐network ansätze, integrate zero‐noise extrapolation and symmetry‐verification error mitigation, and replace ResNet‐18’s initial $7\times7$ convolution with parallel quantum $3\times3$ patch filters. Our approach aims for NISQ‐feasible circuit depths, qubit overhead under 10 per patch, and benchmarks on CIFAR‐100 targeting within 5\% of classical ResNet‐18 accuracy while demonstrating quantum resource savings.
\end{abstract}

\begin{IEEEkeywords}
quantum image representation, NEQR, RGB patch encoding, variational quantum circuits, hybrid ResNet-18
\end{IEEEkeywords}

\section{Background}
High‐quality color image classification critically depends on extracting local color features from small patches. Classical CNNs use $3\times3\times3$ RGB filters to capture spatial and color correlations, a capability missing from current quantum CNNs. In the quantum domain, FRQI encodes grayscale via rotation angles on a single qubit per pixel but loses exact amplitude retrieval~\cite{le2011frqi}. NEQR stores exact $q$‐bit intensities in basis states, requiring $n+q$ qubits for a $2^n\times2^n$ image and scaling linearly with channels~\cite{zhang2013neqr}. Generalized QIR (NGQR) supports arbitrary sizes by adding monitor qubits per dimension but remains single‐channel and further multiplies qubit usage~\cite{xing2024ngqr}. Bit‐plane QIR (BRQI) compresses 8‐bit grayscale into 4 qubits per pixel but cannot preserve arbitrary precision or color channels~\cite{li2018bitplane}. These limitations prevent amplitude‐exact encoding of full RGB patches within NISQ qubit budgets and motivate a 9‐qubit NEQR encoder with compression and trainable filters for hybrid QCNN architectures.

\section{Literature Review}
\subsection{Quantum Image Representations (QIR)}
Early QIR schemes focused on single‐channel images. The Flexible Representation of Quantum Images (FRQI) embeds gray‐scale intensities $\theta_{y,x}$ into rotation angles on a single qubit per pixel, with $n$ position qubits for a $2^n\times2^n$ image~\cite{le2011frqi}. However, FRQI sacrifices exact amplitude retrieval and does not generalize naturally to RGB. The Novel Enhanced QIR (NEQR) encodes each $q$‐bit pixel value exactly into computational basis states across $q$ color qubits plus $2n$ position qubits, but its qubit cost grows linearly with channels and image dimensions~\cite{zhang2013neqr}. To handle arbitrary sizes, the Generalized QIR (NGQR) adds 2–3 monitor qubits per dimension for non‐power‐of‐two images, yet remains single‐channel and still multiplies qubit usage per channel~\cite{xing2024ngqr}. Bit‐plane QIR (BRQI) compacts 8‐bit grayscale into 4 qubits per pixel by bit‐plane slicing, but cannot preserve arbitrary precision or direct color channels~\cite{li2018bitplane}. These limitations highlight the need for a multi‐qubit patch encoder that can amplitude‐exactly load full RGB patches within NISQ budgets.

\subsection{Patch‐Level Quantum Compression}
To reduce qubit load per patch, quantum Principal Component Analysis (qPCA) offers exponential speedup in extracting dominant variance directions, compressing data from 27 values to $k\ll 27$ principal amplitudes in $O(k\,\mathrm{poly}(\log N))$ time~\cite{yu2019quantum}. Variational Quantum Autoencoders (VQAE) further demonstrate end‐to‐end compression and reconstruction of $3\times3$ RGB patches into 5–7 qubits with $>95\%$ fidelity under shallow ansätze~\cite{romero2017quantum}. These patch‐level methods motivate integrating compression with amplitude encoding to lower state‐preparation complexity.

\subsection{Amplitude‐Loading Circuits}
Loading arbitrary amplitudes into an $n$‐qubit state naively requires $O(2^n\,n)$ gates. Gray‐code tree circuits reduce that to $O(n\,2^n)$ depth by hierarchical loading of sparse amplitudes~\cite{plessch2011quantum}. Recent sparse state preparation optimizations eliminate gates for zero components, achieving up to 40\% depth reduction when patches exhibit color redundancy~\cite{zalka1999grover}. Embedding these techniques into NEQR circuits will shrink depth for 27‐value RGB patches.

\subsection{Variational Filter Design and Expressivity}
Quantum Convolutional Neural Networks (QCNN) adapt classical CNN concepts, using patch‐based PQC filters. The original QCNN applied fixed Toffoli and CZ circuits on 10–20 qubits for many‐body state classification~\cite{cong2019qcnn}. Variational Quanvolutional Neural Networks (VQNN) introduced $R_Y$–CNOT layers for $2\times2$ grayscale patches up to 4 qubits, showing minor MNIST accuracy gains but limited expressivity~\cite{henderson2019quanvolutional,mattern2021variational}. Effective rank metrics quantify ansatz expressivity; tree‐tensor‐network (TTN) ansätze aligned to patch connectivity outperform hardware‐efficient layouts by capturing spatial correlations with fewer parameters~\cite{yao2025learning,sugisaki2019quantum}. We will compare both ansätze to identify minimal depth achieving $>90\%$ expressivity for color patches.

\subsection{Error Mitigation Strategies}
NISQ circuits suffer noise‐induced bias. Zero‐noise extrapolation runs each patch circuit at scaled noise levels $\lambda\in\{1,2\}$ and Richardson‐extrapolates $\lambda\to0$ to recover ideal expectation values~\cite{zhu2022efficient}. Symmetry verification discards runs violating invariants—e.g.\ per‐channel mean of a patch amplitude—boosting effective fidelity~\cite{mcclean2018barren}. Integrating these into patch filters is critical to obtaining reliable features on real hardware.

\subsection{Optimizer and Training Dynamics}
Hybrid QNN training faces barren plateaus with increasing qubit count~\cite{mcclean2018barren}. Advanced optimizers like AdaBelief, Yogi, and RMSProp demonstrate smoother convergence in QCNN contexts than Adam or SGD, likely due to adaptive step sizing that counters noise‐driven gradient variance~\cite{zhuang2020adabelief,ceschini2025quanvnn}. We will profile quantum vs.\ classical layer gradients, applying separate learning‐rate schedules to avoid qubit‐scale barren regions.

\subsection{NISQ Resource Mapping}
Mapping patch filters to current devices requires parallel patch‐circuit execution. A $27$‐qubit device can run three $9$‐qubit patch circuits per cycle; a $54$‐qubit device runs six in parallel~\cite{henderson2019quanvolutional,cong2019qcnn}. For a $32\times32$ image with stride 1, $30\times30=900$ patches, processed in batches of 6 yields 150 cycles. With a per‐cycle duration $\sim200\ \mu$s, total filter time $\sim30\,$ms—compatible with ResNet inference latencies.

\subsection{Identified Gaps}
Despite substantial advances in quantum image representations, several critical limitations remain. First, no existing representation encodes exact amplitudes for full 3×3×3 RGB patches within nine qubits or fewer, leaving a gap in multi-channel patch encoding under tight NISQ budgets. Second, many QIR schemes suffer either excessive qubit costs or loss of precision as the number of channels grows, undermining amplitude-exact loading. Third, state-of-the-art QCNN architectures support only grayscale filters of at most four qubits, with no direct color support, preventing color feature extraction at the quantum layer. Fourth, existing variational and fixed-circuit filters lack systematic tools for trainability analysis and expressivity diagnostics, limiting performance optimization. Fifth, no prior work has integrated patch-level quantum filters into modern ResNet-style 3×3 convolution layers, leaving a disconnect between quantum encoders and classical network architectures. Finally, there is a dearth of comprehensive error-mitigation strategies and detailed resource analyses tailored to NISQ hardware requirements, which is essential for practical deployment of quantum filters.

\section{Proposed Direction}
\subsection{Core Innovations}
To address these gaps, we propose a sequence of core innovations. We introduce a nine-qubit NEQR encoder capable of amplitude-exact loading of 3×3×3 RGB patches by leveraging Gray-code amplitude-loading circuits, reducing state-preparation depth by up to forty percent compared with prior tree-decomposition approaches. We integrate patch-level compression via quantum principal component analysis or shallow variational quantum autoencoders, which compress twenty-seven patch amplitudes into six or seven qubits while maintaining fidelity above ninety-five percent. We design trainable variational convolutional filters employing both hardware-efficient and tree-tensor-network ansätze, systematically comparing their expressivity to identify the minimal circuit depth required to exceed ninety-percent effective rank. To mitigate noise on NISQ devices, we combine zero-noise extrapolation over two noise scales with symmetry-verification postselection on per-channel mean invariants. Finally, we embed these quantum filters into a hybrid ResNet-18 architecture by replacing the initial 7×7 convolution and pooling with parallel quantum-encoded 3×3 patch filtering, enabling direct benchmarking against classical ResNet-18 on the CIFAR-100 dataset.

\section{Methodology}
\subsection{Amplitude Encoding}
Load each patch’s 27 intensities $v_i$ into 9 qubits via Gray‐code circuits in $O(9\,2^9)$ gates~\cite{plessch2011quantum}. Apply quantum PCA to compress to 6 qubits when appropriate~\cite{yu2019quantum}.

\subsection{Filter Ansatz}
Design $L\in\{2,3\}$ layers of:
\[
U(\boldsymbol\theta)=\prod_{\ell=1}^{L}
\bigl(\otimes_{q=1}^9R_Y(\theta_{\ell,q})\bigr)
\bigl(\prod_{(q,r)\in E}\mathrm{CNOT}_{q,r}\bigr),
\]
with $E$ matching patch geometry~\cite{sugisaki2019quantum}.

\subsection{Error Mitigation}
Run patch circuits at noise scales $\lambda\in\{1,2\}$, measure expectation $\langle Z\rangle_\lambda$, then extrapolate $\lambda\to0$~\cite{zhu2022efficient}. Enforce mean‐preservation symmetry per channel to postselect valid outcomes.

\subsection{Hybrid ResNet-18}
Our hybrid architecture modifies the standard ResNet-18 by replacing the initial conv$(7 \times 7, s=2)$ plus pooling layer with a parallel quantum patch processing pipeline. The replacement procedure begins by extracting overlapping $3 \times 3$ patches from the input image with unit stride. Each extracted patch undergoes quantum encoding using our 9-qubit NEQR amplitude loading scheme, followed by variational filtering through the designed ansätze, and measurement to obtain classical feature values. The processed patches are then reassembled into $56 \times 56 \times 9$ feature maps that maintain spatial correspondence with the original convolutional output dimensions. Finally, these quantum-generated feature maps proceed through the remaining ResNet-18 blocks without modification, enabling direct comparison with the classical baseline while demonstrating the quantum enhancement at the initial feature extraction stage.

\subsection{Training and Evaluation}
\textbf{Dataset}: CIFAR-100 (50k train, 10k test).  
\textbf{Metrics}: Top-1 accuracy, convergence speed, qubit \& gate counts, latency.  
\textbf{Baselines}: Classical ResNet-18, 4-qubit QCNN, compressed‐embedding QCNN.  
\textbf{Optimizers}: AdaBelief, Yogi, AdaGrad, Adam, RMSProp~\cite{zhuang2020adabelief,ceschini2025quanvnn}.

\section{Expected Contributions}
This research will deliver the first exact nine-qubit quantum encoder for full-color 3×3×3 RGB patches in convolutional neural networks, significantly expanding QCNN capabilities. It will demonstrate efficient patch‐level compression using quantum PCA and variational quantum autoencoders, establishing a new standard for fidelity-preserving qubit reduction. By conducting a comparative expressivity study of hardware-efficient versus tree-tensor-network variational filters, the work will provide actionable guidelines for ansatz selection in quantum image processing. The integration of zero-noise extrapolation and symmetry verification into the quantum filter pipeline will advance error-mitigation practices tailored to variational circuits. Finally, the hybrid quantum-classical ResNet-18 implementation and CIFAR-100 benchmarks will empirically validate NISQ-feasible accuracy within five percent of the classical baseline, demonstrating tangible quantum resource savings.

\bibliographystyle{IEEEtran}
\begin{thebibliography}{99}

\bibitem{le2011frqi}
P. Q. Le, F. Dong, and K. Hirota,
“A flexible representation of quantum images for polynomial preparation, image compression, and processing operations,”
\textit{Quantum Inf. Process.}, vol. 10, pp. 63–84, 2011.

\bibitem{zhang2013neqr}
Y. Zhang, K. Lu, Y. Gao, and M. Wang,
“NEQR: A novel enhanced quantum representation of digital images,”
\textit{Quantum Inf. Process.}, vol. 12, pp. 2833–2860, 2013.

\bibitem{xing2024ngqr}
Z. Xing \textit{et al.},
“NGQR: A novel generalized quantum image representation,”
\textit{IEEE Trans. Emerg. Topics Comput.}, 2024,
doi:10.1109/TETC.2024.3471086.

\bibitem{henderson2019quanvolutional}
M. Henderson, S. Shakya, S. Pradhan, and T. Cook,
“Quanvolutional neural networks: Powering image recognition with quantum circuits,”
\textit{Quantum Mach. Intell.}, vol. 2, no. 1, pp. 1–9, 2020.

\bibitem{mattern2021variational}
D. Mattern \textit{et al.},
“Variational quanvolutional neural networks with enhanced image encoding,”
arXiv:2106.07327, 2021.

\bibitem{fan2024hybrid}
F. Fan \textit{et al.},
“Hybrid quantum-classical convolutional neural network model for image classification,”
\textit{IEEE Trans. Neural Netw. Learn. Syst.}, vol. 35, no. 12, pp. 18145–18157, 2024.

\bibitem{yu2019quantum}
Q. Yu \textit{et al.},
“Quantum data compression by principal component analysis,”
\textit{Quantum Inf. Process.}, vol. 18, p. 045001, 2019.

\bibitem{romero2017quantum}
J. Romero, J. Olson, and A. Aspuru-Guzik,
“Quantum autoencoders for efficient compression of quantum data,”
\textit{Quantum Sci. Technol.}, vol. 2, p. 045001, 2017.

\bibitem{plessch2011quantum}
M. Plesch and \v{C}. Brukner,
“Quantum‐state preparation with universal gate decompositions,”
\textit{Phys. Rev. A}, vol. 83, p. 032302, 2011.

\bibitem{zhu2022efficient}
H. Zhu and H. Zhang,
“Efficient verification of quantum gates with local operations,”
\textit{Phys. Rev. A}, vol. 101, p. 042316, 2020.

\bibitem{mcclean2018barren}
J. R. McClean \textit{et al.},
“Barren plateaus in quantum neural network training landscapes,”
\textit{Nat. Commun.}, vol. 9, p. 4812, 2018.

\bibitem{sugisaki2019quantum}
K. Sugisaki \textit{et al.},
“A probabilistic cloning machine for replicating two non-orthogonal states,”
\textit{Phys. Lett. A}, vol. 243, pp. 261–264, 1998.

\bibitem{zhuang2020adabelief}
H. Zhuang \textit{et al.},
“AdaBelief optimizer: Adapting stepsizes by the belief in observed gradients,”
in \textit{NeurIPS}, 2020.

\bibitem{ceschini2025quanvnn}
A. Ceschini \textit{et al.},
“On hybrid quanvolutional neural networks optimization,”
\textit{Quantum Mach. Intell.}, vol. 7, no. 2, pp. 123–139, 2025.

\bibitem{cong2019qcnn}
I. Cong, S. Choi, and M. D. Lukin,
“Quantum convolutional neural networks,”
\textit{Nat. Phys.}, vol. 15, pp. 1273–1278, 2019.

\bibitem{zalka1999grover}
C. Zalka,
“Grover’s quantum searching algorithm is optimal,”
\textit{Phys. Rev. A}, vol. 60, no. 4, pp. 2746–2751, 1999.

\bibitem{guo2021efficient}
S.-F. Guo \textit{et al.},
“Faster state preparation across quantum phase transition assisted by reinforcement learning,”
\textit{Phys. Rev. Lett.}, vol. 126, p. 060401, 2021.

\bibitem{yao2025learning}
J. Yao,
“Learning to maximize quantum neural network expressivity via effective rank,”
arXiv:2506.15375, 2025.

\bibitem{li2018bitplane}
A.~C. Iliyasu, F.~Yan, and S.~E.~Venegas-Andraca,
“An optimized quantum representation for color digital images,”
\textit{Int. J. Theor. Phys.}, vol. 57, no. 10, pp. 2938–2948, 2018.

\end{thebibliography}

\end{document}
